\documentclass[10pt,letterpaper]{article}
\usepackage[top=2.2cm,bottom=2.2cm,left=6.35cm,right=1.5cm,headheight=14pt,headsep=18pt]{geometry}
\usepackage[utf8]{inputenc}
\usepackage[vietnamese]{babel}
\usepackage{amsmath,amssymb,amsfonts}
\usepackage{xcolor}
\usepackage{tcolorbox}
\usepackage{fancyhdr}
\usepackage{newtxtext,newtxmath}

% Định nghĩa màu chuẩn Stewart Calculus
\definecolor{stewartcyan}{RGB}{0, 160, 220}
\definecolor{stewartred}{RGB}{204, 34, 41}
\definecolor{stewartdarkgray}{RGB}{60, 60, 60}

% Thẻ số phương trình Stewart (hộp viền đỏ, số đỏ không bo góc)
\newcommand{\stewarteqtag}[1]{%
  \tcbox[
    colframe=stewartred,
    colback=white,
    boxrule=0.8pt,
    arc=0pt,
    boxsep=0pt,
    left=3pt,
    right=3pt,
    top=1.5pt,
    bottom=1.5pt
  ]{\fontsize{8.5}{9.5}\selectfont\bfseries\color{stewartred}#1}%
}

% Cấu hình Header và Footer chuẩn trang lẻ sách Stewart
\pagestyle{fancy}
\fancyhf{}
\renewcommand{\headrulewidth}{0pt}
\fancyhead[R]{\textsf{\textbf{MỤC 3.4}}\hspace{1.5em}\textsf{Quy tắc chuỗi}\hspace{2.5em}\textsf{\textbf{205}}}

\setlength{\parindent}{0pt}
\setlength{\parskip}{0pt}

\begin{document}

\noindent
(b) Hàm ngoài là một hàm mũ và hàm trong là hàm bình phương, vì vậy ta sử dụng Công thức~5 và Quy tắc chuỗi để có
\[
h'(x) = \frac{d}{dx}\left(5^{x^2}\right) = 5^{x^2} \ln 5 \cdot \frac{d}{dx}\left(x^2\right) = 2x \cdot 5^{x^2} \ln 5 \tag*{\color{stewartcyan}\rule{1.1ex}{1.1ex}}
\]

\vspace{1.2em}
\noindent
{\color{stewartcyan}\raisebox{0.1ex}{\rule{1.1ex}{1.1ex}}}\hspace{0.5em}\textbf{\textsf{\large Cách chứng minh Quy tắc chuỗi}}

\vspace{0.6em}
\noindent
Nhắc lại rằng nếu $y = f(x)$ và $x$ biến thiên từ $a$ đến $a + \Delta x$, ta định nghĩa số gia của $y$ là
\[
\Delta y = f(a + \Delta x) - f(a)
\]
Theo định nghĩa của đạo hàm, ta có
\[
\lim_{\Delta x \to 0} \frac{\Delta y}{\Delta x} = f'(a)
\]
Do đó, nếu ta ký hiệu $\varepsilon$ là hiệu giữa $\Delta y/\Delta x$ và $f'(a)$, ta thu được
\[
\lim_{\Delta x \to 0} \varepsilon = \lim_{\Delta x \to 0} \left( \frac{\Delta y}{\Delta x} - f'(a) \right) = f'(a) - f'(a) = 0
\]
\noindent
Nhưng\hfill $\varepsilon = \dfrac{\Delta y}{\Delta x} - f'(a) \quad \Rightarrow \quad \Delta y = f'(a)\,\Delta x + \varepsilon\,\Delta x$\hfill\null

\vspace{0.6em}
\noindent
Nếu ta định nghĩa $\varepsilon = 0$ khi $\Delta x = 0$, thì $\varepsilon$ trở thành một hàm liên tục của $\Delta x$. Do đó, đối với một hàm khả vi $f$, ta có thể viết

\vspace{0.4em}
\noindent
\makebox[\linewidth][s]{%
  \stewarteqtag{6}\hfill
  $\displaystyle \Delta y = f'(a)\,\Delta x + \varepsilon\,\Delta x \qquad\text{trong đó}\qquad \varepsilon \to 0 \text{ khi } \Delta x \to 0$\hfill
  \phantom{\stewarteqtag{6}}%
}
\vspace{0.4em}

\noindent
và $\varepsilon$ là một hàm liên tục của $\Delta x$. Tính chất này của các hàm khả vi chính là cơ sở cho phép chúng ta chứng minh Quy tắc chuỗi.

\vspace{1.2em}
\noindent
\textbf{\textsf{\color{stewartred}CHỨNG MINH QUY TẮC CHUỖI}}\quad
Giả sử $u = g(x)$ khả vi tại $a$ và $y = f(u)$ khả vi tại $b = g(a)$. Nếu $\Delta x$ là số gia của $x$, đồng thời $\Delta u$ và $\Delta y$ là các số gia tương ứng của $u$ và $y$, thì ta có thể dùng Phương trình~6 để viết

\vspace{0.4em}
\noindent
\makebox[\linewidth][s]{%
  \stewarteqtag{7}\hfill
  $\displaystyle \Delta u = g'(a)\,\Delta x + \varepsilon_1\,\Delta x = [g'(a) + \varepsilon_1]\,\Delta x$\hfill
  \phantom{\stewarteqtag{7}}%
}
\vspace{0.4em}

\noindent
trong đó $\varepsilon_1 \to 0$ khi $\Delta x \to 0$. Tương tự,

\vspace{0.4em}
\noindent
\makebox[\linewidth][s]{%
  \stewarteqtag{8}\hfill
  $\displaystyle \Delta y = f'(b)\,\Delta u + \varepsilon_2\,\Delta u = [f'(b) + \varepsilon_2]\,\Delta u$\hfill
  \phantom{\stewarteqtag{8}}%
}
\vspace{0.4em}

\noindent
trong đó $\varepsilon_2 \to 0$ khi $\Delta u \to 0$. Bây giờ, nếu thế biểu thức của $\Delta u$ từ Phương trình~7 vào Phương trình~8, ta nhận được
\[
\Delta y = [f'(b) + \varepsilon_2][g'(a) + \varepsilon_1]\,\Delta x
\]
\noindent
do đó\hfill $\dfrac{\Delta y}{\Delta x} = [f'(b) + \varepsilon_2][g'(a) + \varepsilon_1]$\hfill\null

\vspace{0.6em}
\noindent
Khi $\Delta x \to 0$, Phương trình~7 cho thấy $\Delta u \to 0$. Lấy giới hạn khi $\Delta x \to 0$,
\[
\begin{aligned}
\frac{dy}{dx} = \lim_{\Delta x \to 0} \frac{\Delta y}{\Delta x} &= \lim_{\Delta x \to 0}\,[f'(b) + \varepsilon_2][g'(a) + \varepsilon_1] \\[5pt]
&= f'(b)\,g'(a) = f'(g(a))\,g'(a)
\end{aligned}
\]
\noindent
Điều này chứng minh Quy tắc chuỗi.\hfill{\color{stewartred}\raisebox{0.1ex}{\rule{1.1ex}{1.1ex}}}

\end{document}